\subsection{Compilation flow}\label{subsec:compilation}

Every benchmark circuit is built once by \texttt{bench/circuits.py} (fixed seeds) and passed unchanged to the FPGA
engines and to the Qiskit CPU runs of Section~\ref{subsec:res-cpu}; QASM copies are part of the reproduction package.
The flow from a Qiskit circuit to engine records has five steps, all on the host and all included in
$T_{\mathrm{compile}}$ (Table~\ref{tab:timing}):
\begin{enumerate}
\item \emph{Gate intake.} Gates of one or two qubits are taken as unitary matrices from Qiskit
(\texttt{Operator}); multi controlled $X$, $Z$ and phase gates are taken directly as one controlled $2\times2$
operation. Such a gate stays one record as long as its target and its local controls fit the tile window and its
other controls sit on slice qubits, where the slice enable (F2) applies them; with $k$ tile qubits and $s$ slice qubits
this allows gates of up to $k+s$ qubits. Wider gates, every gate wider than the tile window when F2 is switched off,
and gates without a controlled form are decomposed by the Qiskit transpiler with \texttt{basis\_gates=[u, cx]},
\texttt{optimization\_level=1}. No other transpilation is applied, so the
FPGA and CPU runs see the same gate sequence.
\item \emph{Fusion and relabelling.} Runs of single qubit gates on one qubit are multiplied into one matrix; SWAP
gates are absorbed into the logical to physical qubit map and cost nothing.
\item \emph{Classification.} With the current qubit map each operation is classified as local, slice controlled
(F2), slice diagonal (F1), slice permutation (F3) or exchange (Section~\ref{subsec:engine}).
\item \emph{Scheduling.} Operations are appended to the current pass while their qubits fit the tile window;
an exchange operation first triggers a remap that brings its slice qubit to a local position, choosing as victims the
local qubits used furthest in the future, and the remapped qubit stays local for all later gates on it. A multi controlled gate wider than the tile window is placed by one remap that swaps its target with one of
its own controls, which then acts through the slice enable, and moves surplus local controls to free slice positions;
this keeps the Grover oracle a single record up to $n=14$ on one and two cards and $n=16$ on four cards.
\item \emph{Encoding.} Passes between remaps form one segment, uploaded once to every engine and executed by one
launch per engine.
\end{enumerate}

In a statevector simulator the cost is set by the number of full sweeps over the state and by the amplitudes that
cross slices, not by circuit depth. GHZ preparation, a serial chain of CX gates, shows what this means for the
compilation. On the engine all CX gates of the textbook GHZ chain whose qubits are local
share one pass, so depth does not cost sweeps; what costs is each CX whose target is a slice qubit, because the
target must first be brought into the slice (one remap per slice qubit, \NumGhzChainRemapsNoFfour~remaps with four
cards and \NumEnginesFlagship~engines per card). A logarithmic depth construction~\cite{cruz2019efficient} reduces depth but
leaves this count almost unchanged (\NumGhzLogRemapsNoFfour). The compiler therefore applies a hardware specific
rewrite (F4): for an $X$ type gate (CX, Toffoli) whose target is a slice qubit and whose controls are local, it moves
the controls up into slice positions instead of moving the target down. The remap costs the same, but the gate and
every later CX among these qubits then become slice relabelling (F3) or slice controlled (F2) operations; the GHZ
chain needs \NumGhzChainRemapsFfour~remap at every width. Because the rewrite can also increase traffic (it does for
the controlled modular multipliers of Shor's circuit), the compiler schedules every circuit with and without F4 and
keeps the schedule with fewer remaps. Section~\ref{subsec:res-ablation} reports the measured times of all variants.
